\section{September 2, 2026}

We develop the coefficient representation of Runge--Kutta methods and
derive their order conditions. Throughout, we retain the notation of
the previous lectures: \(\tau_n=t_{n+1}-t_n\), \(x_n\) is the numerical
approximation at \(t_n\), \(\Phi\) is the exact flow, and \(\Psi\) is
the numerical one-step map.

\subsection{Butcher tableaux}

An \(s\)-stage Runge--Kutta method has the form
\begin{align}
  k_i
  &=f\left(t_n+c_i\tau_n,
    x_n+\tau_n\sum_{j=1}^s a_{ij}k_j\right),
    \qquad i=1,\ldots,s,
    \label{eq:rk-stages-sep-02}\\
  x_{n+1}
  &=x_n+\tau_n\sum_{i=1}^s b_i k_i.
    \label{eq:rk-update-sep-02}
\end{align}
The coefficients are collected into
\[
  A=(a_{ij})_{i,j=1}^s,
  \qquad
  b=\begin{bmatrix}b_1&\cdots&b_s\end{bmatrix}^{\mathsf T},
  \qquad
  c=\begin{bmatrix}c_1&\cdots&c_s\end{bmatrix}^{\mathsf T}.
\]
They are conventionally displayed in a \emph{Butcher tableau}:
\[
  \begin{array}{c|cccc}
    c_1&a_{11}&a_{12}&\cdots&a_{1s}\\
    c_2&a_{21}&a_{22}&\cdots&a_{2s}\\
    \vdots&\vdots&\vdots&\ddots&\vdots\\
    c_s&a_{s1}&a_{s2}&\cdots&a_{ss}\\ \hline
       &b_1&b_2&\cdots&b_s
  \end{array}.
\]
The values \(t_n+c_i\tau_n\) are the intermediate stage times between
\(t_n\) and \(t_{n+1}\). If \(A\) is strictly lower triangular, then
each stage depends only on previously computed stages. The resulting
method is an \emph{explicit Runge--Kutta method}.

\subsection{The internal consistency relation}

Any nonautonomous ODE can be converted into an autonomous system by
augmenting its state. Set
\[
  z(t)=\begin{bmatrix}t\\x(t)\end{bmatrix},
  \qquad
  z'(t)=
  \begin{bmatrix}1\\f(t,x(t))\end{bmatrix}
  =F(z(t)).
\]
Apply a Runge--Kutta method to this autonomous system. Because the time
component of every stage slope is one, the time component of stage
\(i\) is
\[
  t_n+\tau_n\sum_{j=1}^s a_{ij}.
\]
For this to agree with the prescribed stage time
\(t_n+c_i\tau_n\) in \eqref{eq:rk-stages-sep-02}, we require
\[
  c_i=\sum_{j=1}^s a_{ij},
  \qquad i=1,\ldots,s.
\]
Thus the \emph{internal consistency relation} is
\begin{equation}\label{eq:internal-consistency-sep-02}
  c=A\mathbf{1},
  \qquad
  \mathbf{1}=\begin{bmatrix}1&\cdots&1\end{bmatrix}^{\mathsf T}.
\end{equation}
Under this relation, applying the Runge--Kutta scheme directly to the
nonautonomous equation is equivalent to applying it to the augmented
autonomous system.

\subsection{One-step error and order}

For an exact starting value \(x\) at time \(t\), define the one-step
defect by
\begin{equation}\label{eq:one-step-defect-sep-02}
  \mathcal{E}(x,t,\tau)
  =\Phi^{t+\tau,t}(x)-\Psi^{t+\tau,t}(x).
\end{equation}
A Runge--Kutta method has order \(p\) if
\[
  \mathcal{E}(x,t,\tau)=O(\tau^{p+1})
  \qquad\text{as }\tau\to0.
\]
Assuming \eqref{eq:internal-consistency-sep-02}, it suffices to derive
the conditions for an autonomous equation
\[
  x'=f(x).
\]
We write its exact and Runge--Kutta flows as \(\Phi^\tau\) and
\(\Psi^\tau\), respectively.

Let \(f'=Df(x)\) and \(f''=D^2f(x)\). Differentiation along an exact
solution gives
\begin{align*}
  x'&=f,\\
  x''&=f'f,\\
  x'''&=f''(f,f)+f'f'f.
\end{align*}
Consequently,
\begin{equation}\label{eq:exact-flow-third-order-sep-02}
  \Phi^\tau(x)
  =x+\tau f+\frac{\tau^2}{2}f'f
   +\frac{\tau^3}{6}\bigl(f''(f,f)+f'f'f\bigr)
   +O(\tau^4).
\end{equation}

For the numerical method, set
\[
  \Delta_i=\tau\sum_{j=1}^s a_{ij}k_j,
  \qquad
  k_i=f(x+\Delta_i).
\]
Since \(k_j=f+O(\tau)\) and \(c=A\mathbf{1}\),
\[
  \Delta_i=\tau c_i f+O(\tau^2).
\]
A more detailed Taylor expansion gives
\[
  k_i
  =f+\tau c_i f'f
   +\tau^2\left(
     \sum_{j=1}^s a_{ij}c_j\,f'f'f
     +\frac{c_i^2}{2}f''(f,f)
   \right)
   +O(\tau^3).
\]
Substitution into \(\Psi^\tau(x)=x+\tau\sum_i b_i k_i\) yields
\begin{align}
  \Psi^\tau(x)
  ={}&x+\tau(b^{\mathsf T}\mathbf{1})f
  +\tau^2(b^{\mathsf T}c)f'f\notag\\
  &+\tau^3\left(
    (b^{\mathsf T}Ac)f'f'f
    +\frac12(b^{\mathsf T}c^{\circ2})f''(f,f)
  \right)
  +O(\tau^4),
  \label{eq:rk-flow-third-order-sep-02}
\end{align}
where \(c^{\circ2}=(c_1^2,\ldots,c_s^2)^{\mathsf T}\).

Comparing \eqref{eq:exact-flow-third-order-sep-02} and
\eqref{eq:rk-flow-third-order-sep-02} gives the order conditions through
order three:
\begin{alignat}{3}
  b^{\mathsf T}\mathbf{1}&=1,
  &\qquad&\text{(order one)},
  \label{eq:rk-oc1-sep-02}\\
  b^{\mathsf T}c&=\frac12,
  &&\text{(additional condition for order two)},
  \label{eq:rk-oc2-sep-02}\\
  b^{\mathsf T}c^{\circ2}&=\frac13,
  &&\text{(additional condition for order three)},
  \label{eq:rk-oc3-sep-02}\\
  b^{\mathsf T}Ac&=\frac16,
  &&\text{(additional condition for order three)}.
  \label{eq:rk-oc4-sep-02}
\end{alignat}
Therefore \eqref{eq:rk-oc1-sep-02} gives a first-order method; the first
two conditions give a second-order method; and all four give a
third-order method. A fourth-order method must satisfy eight conditions
in total. For nonautonomous equations, the relation
\(c=A\mathbf{1}\) is also needed for these autonomous order conditions
to apply directly.

\subsection{Examples}

\begin{example}[Forward Euler]
  Forward Euler has the tableau
  \[
    \begin{array}{c|c}
      0&0\\ \hline
       &1
    \end{array}.
  \]
  Thus \(s=1\), \(A=0\), \(b=1\), and \(c=0\). It satisfies
  \(b^{\mathsf T}\mathbf{1}=1\), but not the second-order condition,
  so it is first order.
\end{example}

\begin{example}[Explicit trapezoidal method]
  The explicit trapezoidal method, or Heun's method, is
  \begin{align*}
    k_1&=f(t_n,x_n),\\
    k_2&=f(t_n+\tau_n,x_n+\tau_n k_1),\\
    x_{n+1}&=x_n+\frac{\tau_n}{2}(k_1+k_2).
  \end{align*}
  Its tableau is
  \[
    \begin{array}{c|cc}
      0&0&0\\
      1&1&0\\ \hline
       &\tfrac12&\tfrac12
    \end{array}.
  \]
  Here \(c=A\mathbf{1}\),
  \(b^{\mathsf T}\mathbf{1}=1\), and
  \(b^{\mathsf T}c=1/2\). However,
  \(b^{\mathsf T}c^{\circ2}=1/2\neq1/3\), so the method is second
  order but not third order.
\end{example}

\subsection{Explicit three-stage methods}

A general explicit three-stage Runge--Kutta method has tableau
\[
  \begin{array}{c|ccc}
    0   &0&0&0\\
    c_2 &a_{21}&0&0\\
    c_3 &a_{31}&a_{32}&0\\ \hline
        &b_1&b_2&b_3
  \end{array}.
\]
Internal consistency requires
\[
  c_2=a_{21},
  \qquad
  c_3=a_{31}+a_{32}.
\]
Before imposing the order conditions, the six free coefficients are
\[
  b_1,b_2,b_3,a_{21},a_{31},a_{32}.
\]
The four conditions
\eqref{eq:rk-oc1-sep-02}--\eqref{eq:rk-oc4-sep-02} enforce third order,
so one generically expects a two-parameter family of explicit
three-stage schemes of order three.
